\section{Introduction}
\label{sec:introduction}

Let $\pi:X\to Y$ be a factor map between one-dimensional shifts of finite
type.  We call it \emph{split} if it admits a continuous shift-commuting
section $\iota:Y\to X$ with $\pi\circ\iota=\id_Y$.  A section makes $Y$
visible as a copy inside $X$, but by itself gives little quantitative
information about the entropy excess
\[
 \Gap(\pi)=h(X)-h(Y).
\]
This note asks for a local split construction on a fixed presentation for
which the finite equation and the leading entropy constant can both be read
off explicitly.

The construction is elementary to state.  Let $G$ be a finite directed
multigraph with primitive adjacency matrix $A$, fix a vertex $v$, and choose
an old length-$N$ return walk $\gamma_N$ based at $v$.  Adjoin an internally
new return chain of exactly $N$ edges based at $v$, with all $N-1$ internal
vertices distinct and disjoint from $G$.  Map the $j$th new edge to the $j$th edge of
$\gamma_N$ and fix the old graph.  Equal edge lengths make this an
incidence-preserving graph morphism.  The induced one-block factor is split
by old-graph inclusion and is noninjective because the old and new periodic
walks have the same image.

The exact finite equation is most transparent at the entropy radius.  Put
\[
 D(z)=\det(I-zA),\qquad
 C_v(z)=\det(I-zA^{(v)}),
\]
where $A^{(v)}$ deletes row and column $v$.  If $A_N$ is the enlarged
adjacency matrix, then
\begin{equation}
 \det(I-zA_N)=D(z)-z^NC_v(z).
 \label{eq:intro-determinant}
\end{equation}
Thus the sign is minus and the exponent counts added path \emph{edges}, not
internal vertices.  If $F_v$ is the old first-return generating function and
$r_N=\rho(A_N)^{-1}$, the equivalent renewal equation is
\begin{equation}
 F_v(r_N)+r_N^N=1.
 \label{eq:intro-renewal}
\end{equation}
The renewal determinant mechanism itself is classical; related determinant
ratios and first-return or \emph{rome} methods appear in
\cite{Gomez2010,BlockEtAl1980}.  We prove the identities because their sign,
exponent, and interpretation are essential to the quantitative result.

Suppose first that $\lambda=\rho(A)>1$ and put $r=\lambda^{-1}$.  For
positive right and left Perron eigenvectors $x,y$, define
\begin{equation}
 \pi_v=\frac{x_vy_v}{y^Tx}.
 \label{eq:intro-pi}
\end{equation}
The diagonal Perron resolvent residue gives
$F_v'(r)=\lambda/\pi_v$.  We prove the relative-error law
\begin{equation}
 \Gap_N:=h(X_{G_N})-h(X_G)
 =\pi_v\lambda^{-N}\{1+O(N\lambda^{-N})\}.
 \label{eq:intro-positive-law}
\end{equation}
In particular, the ratio of the two sides tends to one.  The use of both
left and right eigenvectors is mandatory for nonsymmetric graphs.

At the mixing zero-entropy boundary, the essential graph is the one-loop
graph.  The exact equation becomes
\begin{equation}
 \ee^{-h_N}+\ee^{-Nh_N}=1.
 \label{eq:intro-zero-equation}
\end{equation}
Writing $w_N=W(N)$ for the principal Lambert function, we obtain
\begin{equation}
 h_N=\frac{w_N}{N}
 +\frac{w_N^2}{2N^2(w_N+1)}
 +O\left(\frac{w_N^2}{N^3}\right).
 \label{eq:intro-zero-law}
\end{equation}
Thus the gap scale changes from exponential to $W(N)/N$.  The estimate
$Nr^N\to0$ used in \cref{eq:intro-positive-law} fails exactly when $r=1$;
the boundary law cannot be obtained by substitution.

\begin{table}[t]
\caption{The local construction is uniform, but its entropy asymptotics split
at the zero-entropy boundary.}
\label{tab:regimes}
\centering
\small
\renewcommand{\arraystretch}{1.12}
\begin{tabularx}{\textwidth}{@{}>{\raggedright\arraybackslash}p{0.18\textwidth}
>{\raggedright\arraybackslash}p{0.23\textwidth}
>{\raggedright\arraybackslash}p{0.27\textwidth}X@{}}
\toprule
Regime & Old essential graph & Gap law & Mandatory firewall \\
\midrule
Positive entropy
& primitive $A$, $\lambda>1$
& $\Gap_N=\pi_v\lambda^{-N}\{1+O(N\lambda^{-N})\}$
& $\pi_v$ uses left and right PF eigenvectors \\
Mixing zero entropy
& one loop, $A=[1]$
& $h_N=W(N)/N$ plus an explicit first correction
& never substitute $r=1$ in the exponential proof \\
\bottomrule
\end{tabularx}
\end{table}



\paragraph{Ownership and residual scope.}
Recent work on factorizable embeddings and retractions includes
\cite{MarcusEtAl2025}; map-extension and contractibility perspectives include
\cite{Meyerovitch2025,PoirierSalo2024}.  These works frame a much broader
structural theory.  We do not give criteria deciding whether an arbitrary
factor admits a section.  Nor is the existence of arbitrarily small split
entropy gaps a residual claim: projection
$Y\times Z_\varepsilon\to Y$ splits whenever $Z_\varepsilon$ has a fixed
point, and subsystem-entropy density is developed in
\cite{BlandMcGoffPavlov2023}.  The candidate contribution is only the
combined presentation-local chain construction,
\cref{eq:intro-determinant,eq:intro-renewal}, and the normalized two-regime
laws \cref{eq:intro-positive-law,eq:intro-zero-law}.  No priority conclusion
is drawn from a search miss; external release remains on hold for specialist
clearance.

\paragraph{Organization.}
\Cref{sec:rigidity} records the supporting entropy-rigidity fact and its
necessary irreducibility hypothesis.  \Cref{sec:construction} defines the
one-block split construction.  \Cref{sec:exact} proves the determinant and
renewal equations.  \Cref{sec:positive} derives the positive-entropy
constant, while \cref{sec:zero} treats the boundary regime.
\Cref{sec:controls} reports deterministic controls and scope limits.
